\documentclass[11pt]{article}

\usepackage[T1]{fontenc}
\usepackage[utf8]{inputenc}
\usepackage{lmodern}
\usepackage{microtype}
\usepackage[margin=1in,headheight=15pt]{geometry}
\usepackage{amsmath,amssymb,mathtools}
\usepackage{booktabs,longtable,tabularx,array}
\usepackage[dvipsnames]{xcolor}
\usepackage{enumitem}
\usepackage{fancyhdr}
\usepackage{hyperref}

\definecolor{navy}{HTML}{17365D}
\definecolor{teal}{HTML}{137C8B}
\definecolor{softblue}{HTML}{EAF1F8}
\definecolor{softgray}{HTML}{F3F5F7}
\definecolor{amber}{HTML}{8A5A00}
\definecolor{green}{HTML}{176B3A}
\definecolor{red}{HTML}{9A2D2D}

\hypersetup{
  colorlinks=true,
  hypertexnames=false,
  linkcolor=navy,
  urlcolor=teal,
  pdftitle={Mortensen--Pissarides (1994): Readable Lean 4 Formalization},
  pdfauthor={Generated from the Lean4Tutorial MP1994 development}
}

\pagestyle{fancy}
\fancyhf{}
\fancyhead[L]{\small Mortensen--Pissarides (1994) in Lean 4}
\fancyhead[R]{\small Mathematical audit}
\fancyfoot[C]{\thepage}

\setlength{\parindent}{0pt}
\setlength{\parskip}{0.55em}
\setlist{nosep,leftmargin=1.5em}
\setcounter{tocdepth}{2}
\allowdisplaybreaks
\sloppy

\newcommand{\code}[1]{\texttt{#1}}
\newcommand{\breakcode}[1]{\nolinkurl{#1}}
\newcommand{\leanfile}[1]{\texttt{#1}}
\newcommand{\epsd}{\varepsilon_d}
\newcommand{\epsu}{\varepsilon_u}
\newcommand{\mworker}[1]{m(#1,1)}
\newcommand{\status}[2]{\textcolor{#1}{\textbf{#2}}}
\newcommand{\decl}[1]{\par\smallskip\noindent\colorbox{softblue}{\parbox{0.965\linewidth}{\textbf{Lean declaration:} \code{#1}}}\par\smallskip}
\newcommand{\source}[1]{\textcolor{navy}{\small\textbf{Source:} \leanfile{#1}}}

\title{\vspace{-2em}\textcolor{navy}{\textbf{Mortensen--Pissarides (1994)}}\\[0.35em]
\Large A Readable Mathematical Rendering of the Lean 4 Formalization}
\author{Generated from \code{Lean4Tutorial.i002\_replicate\_MP}}
\date{20 July 2026}

\begin{document}
\maketitle

\begin{center}
\colorbox{softgray}{\parbox{0.91\linewidth}{
\textbf{Purpose.} This document translates the complete mathematical interface of the
Lean 4 development into conventional notation. It is intended for manual comparison
with Dale T. Mortensen and Christopher A. Pissarides, ``Job Creation and Job
Destruction in the Theory of Unemployment,'' \emph{Review of Economic Studies}
61 (1994), 397--415. It is not a replacement for the Lean kernel check; it is the
human-readable companion to that check.
}}
\end{center}

\vfill
\begin{tabularx}{\linewidth}{@{}>{\bfseries}p{0.27\linewidth}X@{}}
Lean build target & \code{lake build Lean4Tutorial.i002\_replicate\_MP} \\
Kernel-checked results & 31 named theorems, each audited with \code{assert\_no\_sorry} \\
Mathematical source & \leanfile{Definitions.lean}, \leanfile{Equilibrium.lean}, and \leanfile{Results.lean} \\
Verification source & \leanfile{Verification.lean} \\
Document status & \status{green}{Built from the current Lean sources} \\
\end{tabularx}
\vfill

\textbf{Reading warning.} A Lean \emph{structure field} is assumed when an object of
that structure is supplied. It is not automatically a derived theorem. This matters
for the cyclical equilibrium structures, where some comparative-static orderings are
encoded as fields. Section~\ref{sec:scope} separates assumptions, definitions, and
derived results.

\newpage
\tableofcontents
\newpage

\section{How to read this audit}

Every substantive formula below is paired with its Lean declaration. The translation
uses ordinary mathematical notation while preserving the exact content of the Lean
term. In particular:

\begin{itemize}
  \item a \textbf{definition} introduces notation or a proposition;
  \item an \textbf{equilibrium structure} packages variables together with conditions
        that must be provided by any instance;
  \item a \textbf{theorem} is a kernel-checked implication from explicit hypotheses;
  \item an \code{assert\_no\_sorry} check rules out dependence on Lean's
        \code{sorryAx} for the named theorem.
\end{itemize}

The Lean source uses \code{P.«lambda»} internally for the paper's $\lambda$, because
bare $\lambda$ is reserved syntax in Lean. This document always prints the paper
symbol $\lambda$.

\subsection{Source files}

\begin{longtable}{@{}p{0.29\linewidth}p{0.66\linewidth}@{}}
\toprule
\textbf{File} & \textbf{Role} \\
\midrule
\endhead
\leanfile{Definitions.lean} & Primitives, matching technology, Bellman predicates,
reduced equilibrium equations, gross flows, and period accounting. \\
\leanfile{Equilibrium.lean} & Stationary value equilibrium, reduced steady state,
two-state and Markov interfaces, and simulation paths. \\
\leanfile{Results.lean} & All derived propositions currently proved by the development. \\
\leanfile{Verification.lean} & One \code{assert\_no\_sorry} command for each of the
31 theorems in \leanfile{Results.lean}. \\
\leanfile{i002\_replicate\_MP.lean} & Dedicated build entry point importing the four
modules above. \\
\bottomrule
\end{longtable}

\section{Primitives and maintained assumptions}

\source{Definitions.lean: Primitives, Assumptions}

An instance $P$ of \code{Primitives} contains the following exogenous objects.

\begin{longtable}{@{}>{$}p{0.12\linewidth}<{$}p{0.31\linewidth}p{0.48\linewidth}@{}}
\toprule
\textbf{Symbol} & \textbf{Lean field} & \textbf{Meaning} \\
\midrule
\endhead
r & \code{r} & Continuous-time discount rate. \\
\lambda & \code{«lambda»} & Poisson arrival rate of job-specific shocks. \\
\sigma & \code{sigma} & Dispersion of job-specific productivity. \\
\beta & \code{beta} & Worker's Nash share of match surplus. \\
c & \code{c} & Flow cost of maintaining a vacancy. \\
b & \code{b} & Unemployment income or value of leisure. \\
p & \code{p} & Aggregate component of productivity. \\
\epsu & \code{epsilon\_u} & Upper support point of the idiosyncratic shock. \\
dF & \code{dF} & Probability measure generating the shock distribution. \\
q(\theta) & \code{q} & Vacancy meeting rate. \\
\bottomrule
\end{longtable}

The maintained \code{Assumptions P} proposition is exactly
\begin{align*}
&r>0,\qquad \lambda\ge 0,\qquad \sigma>0,\qquad 0<\beta<1,
\qquad c>0,\\
&dF(\mathbb R)=1,\qquad dF(({-\infty},\epsu])=1,\\
&0\le F(x)\le 1,\qquad F\text{ is nondecreasing},\\
&\theta>0\Longrightarrow q(\theta)>0,\qquad q\text{ is nonincreasing},\\
&\theta\longmapsto \theta q(\theta)\text{ is nondecreasing}.
\end{align*}
The CDF bounds and monotonicity are fields, rather than results derived in Lean from
the measure definition. This is a deliberate strengthening of the interface.

\section{Basic mathematical definitions}

\source{Definitions.lean: namespace MP1994.Primitives}

\subsection{Distribution, matching, and productivity}

\begin{align*}
F(x) &:= dF(({-\infty},x]), \\
\theta(v,u) &:= \frac vu, \\
\mworker{\theta} &:= \theta q(\theta), \\
m(v,u) &:= vq(v/u), \\
y(\varepsilon) &:= p+\sigma\varepsilon.
\end{align*}

The Lean names are respectively \code{F}, \code{theta}, \code{mWorker}, \code{m},
and \code{price}. Notice that the code's aggregate matching definition is
$v q(v/u)$, while the worker meeting rate is $\theta q(\theta)$.

\subsection{Continuation values and separations}

\begin{align*}
I(\epsd) &:= \int_{\epsd}^{\epsu}[1-F(x)]\,dx,
&&\text{\code{optionValueIntegral}},\\
T(\ell,h) &:= \int_{\ell}^{h}[1-F(x)]\,dx,
&&\text{\code{tailIntegral}},\\
\mathcal P(S) &:= \int_{\mathbb R}\max\{S(x),0\}\,dF(x),
&&\text{\code{positiveContinuation}},\\
s(\epsd) &:= \lambda F(\epsd),
&&\text{\code{destructionHazard}}.
\end{align*}

The first two integrals are Lebesgue interval integrals with respect to $dx$;
$\mathcal P(S)$ is an integral with respect to the shock measure $dF$.

\subsection{Closed-form cutoff surplus}

\begin{align*}
S_d(\varepsilon;\epsd)
  &:= \frac{\sigma(\varepsilon-\epsd)}{r+\lambda},
  &&\text{\code{S\_cutoff}},\\
J_d(\varepsilon;\epsd)
  &:= (1-\beta)S_d(\varepsilon;\epsd),
  &&\text{\code{J\_cutoff}},\\
[W-U]_d(\varepsilon;\epsd)
  &:= \beta S_d(\varepsilon;\epsd),
  &&\text{\code{WminusU\_cutoff}}.
\end{align*}

\section{Stationary value equations}

\source{Definitions.lean and ValueEquilibrium in Equilibrium.lean}

The structure \code{ValueEquilibrium P} supplies
$\theta,\epsd,V,U,J,W,w,S$, with $\theta>0$ and $\epsd<\epsu$, and requires all
conditions below.

\decl{Primitives.VacancyBellman --- paper equation (1)}
\begin{equation}
rV=-c+q(\theta)\bigl[J(\epsu)-V\bigr]. \tag{1}
\end{equation}

\decl{Primitives.FreeEntry --- paper equation (2)}
\begin{equation}
rV=0. \tag{2}
\end{equation}

\decl{Primitives.SurplusIdentity --- paper equation (3)}
\begin{equation}
S(\varepsilon)=J(\varepsilon)+W(\varepsilon)-U
\qquad\text{for every }\varepsilon. \tag{3}
\end{equation}

\decl{Primitives.NashSharing --- paper equation (4)}
\begin{equation}
W(\varepsilon)-U=\beta S(\varepsilon)
\qquad\text{for every }\varepsilon. \tag{4}
\end{equation}

\decl{Primitives.FilledJobBellman --- paper equation (5)}
\begin{equation}
rJ(\varepsilon)=p+\sigma\varepsilon-w(\varepsilon)
+\lambda(1-\beta)\bigl[\mathcal P(S)-S(\varepsilon)\bigr]. \tag{5}
\end{equation}

\decl{Primitives.WorkerBellman --- paper equation (6)}
\begin{equation}
rW(\varepsilon)=w(\varepsilon)
+\lambda\beta\bigl[\mathcal P(S)-S(\varepsilon)\bigr]. \tag{6}
\end{equation}

\decl{Primitives.UnemployedBellman --- paper equation (7)}
\begin{equation}
rU=b+\mworker{\theta}\,[W(\epsu)-U]. \tag{7}
\end{equation}

\decl{Primitives.SurplusBellman --- paper equation (8)}
\begin{equation}
(r+\lambda)S(\varepsilon)
=p+\sigma\varepsilon-b+\lambda\mathcal P(S)
-\beta\mworker{\theta}S(\epsu). \tag{8}
\end{equation}

The structure additionally requires
\begin{equation*}
S(\epsd)=0,
\qquad
S(\varepsilon)\ge0\ \Longleftrightarrow\ \epsd\le\varepsilon.
\end{equation*}
These are \code{cutoff\_zero} and \code{reservation\_rule}.

\section{Reduced stationary equilibrium}

\source{Definitions.lean and SteadyStateEquilibrium in Equilibrium.lean}

\decl{Primitives.JobDestructionCondition --- paper equation (10)}
\begin{equation}
p+\sigma\epsd
=b+\frac{\beta c}{1-\beta}\theta
-\frac{\sigma\lambda}{r+\lambda}
  \int_{\epsd}^{\epsu}[1-F(x)]\,dx. \tag{10}
\end{equation}

\decl{Primitives.dispersionCutoffResponse --- right side of paper equation (11)}
The code defines the quantity described in its documentation as
$\sigma\,\partial\epsd/\partial\sigma$ at fixed $\theta$:
\begin{equation}
R_\sigma(\theta,\epsd)
:=\frac{(r+\lambda)/\sigma}{r+\lambda F(\epsd)}
\left[p-b-\frac{\beta c}{1-\beta}\theta\right]. \tag{11R}
\end{equation}
Lean defines the right-hand side; it does not define a derivative operator or prove
that an implicit derivative equals this expression.

\decl{Primitives.S\_cutoff --- paper equation (12)}
\begin{equation}
S_d(\varepsilon;\epsd)
=\frac{\sigma(\varepsilon-\epsd)}{r+\lambda}. \tag{12}
\end{equation}

\decl{Primitives.JobCreationCondition --- paper equation (13)}
\begin{equation}
q(\theta)
=\frac{c}{1-\beta}\,
  \frac{r+\lambda}{\sigma(\epsu-\epsd)}. \tag{13}
\end{equation}

\decl{Primitives.BeveridgeCondition --- paper equation (14)}
\begin{equation}
u=\frac{\lambda F(\epsd)}
        {\lambda F(\epsd)+\mworker{\theta}}. \tag{14}
\end{equation}

\decl{Primitives.unemploymentDrift --- paper equation (15)}
\begin{equation}
\dot u=(1-u)\lambda F(\epsd)-u\mworker{\theta}. \tag{15}
\end{equation}

The structure \code{SteadyStateEquilibrium P} is a triple $(\theta,\epsd,u)$
satisfying (10), (13), and (14), together with
$\theta>0$, $\epsd<\epsu$, and $0\le u\le1$.
Its derived notation is
\begin{equation*}
v:=\theta u,\qquad n:=1-u,\qquad
C:=u\mworker{\theta},\qquad
D:=(1-u)\lambda F(\epsd).
\end{equation*}

\section{Anticipated aggregate states and Markov equations}

\source{Definitions.lean and Equilibrium.lean}

Stars denote boom values and unstarred quantities denote recession values.

\decl{Primitives.BoomDestructionCondition --- paper equation (20)}
\begin{align}
p^*+\sigma\epsd^*
={}&b+\frac{\beta c}{1-\beta}\theta^*
-\frac{\lambda\sigma}{r+\lambda+\mu}
  \int_{\epsd^*}^{\epsd}[1-F(x)]\,dx \notag\\
&-\frac{\lambda\sigma}{r+\lambda}
  \int_{\epsd}^{\epsu}[1-F(x)]\,dx. \tag{20}
\end{align}

\decl{Primitives.RecessionDestructionCondition --- paper equation (24)}
\begin{align}
p+\sigma\epsd
={}&b+\frac{\beta c}{1-\beta}\theta
-\frac{\lambda\sigma}{r+\lambda}
  \int_{\epsd}^{\epsu}[1-F(x)]\,dx \notag\\
&-\frac{\mu\sigma}{r+\lambda+\mu}(\epsd-\epsd^*). \tag{24}
\end{align}

\decl{Primitives.boomSurplus --- paper equation (29)}
\begin{equation}
S^*(\varepsilon)
=\frac{\sigma(\varepsilon-\epsd^*)}{r+\lambda+\mu}
+\frac{\mu}{r+\lambda+\mu}S_d(\varepsilon;\epsd). \tag{29}
\end{equation}

\decl{Primitives.BoomJobCreationCondition --- paper equation (30)}
\begin{equation}
q(\theta^*)
=\frac{c(r+\lambda)/(1-\beta)}
{\sigma(\epsu-\epsd^*)
-\dfrac{\mu\sigma(\epsd-\epsd^*)}{r+\lambda+\mu}}. \tag{30}
\end{equation}

\subsection{Two-state structure fields}

\code{AnticipatedTwoStateEquilibrium P} packages $(p,p^*,\mu,\theta,\theta^*,
\epsd,\epsd^*)$ and requires
\begin{gather*}
p<p^*,\quad \mu\ge0,\quad \theta>0,\quad\theta^*>0,
\quad \epsd^*<\epsd<\epsu,\\
\text{equations (20), (24), (13) for recession, and (30) for the boom.}
\end{gather*}
The simpler \code{TwoStateEquilibrium recession boom} packages two reduced steady
states and directly requires
\begin{equation*}
p_{\rm rec}<p_{\rm boom},\qquad
\theta_{\rm rec}<\theta_{\rm boom},\qquad
\varepsilon_{d,{\rm boom}}<\varepsilon_{d,{\rm rec}}.
\end{equation*}
Thus these three orderings are data of that structure, not the conclusions of a
constructor theorem in the current development.

\subsection{General finite-state value equilibrium}

For a finite state type, \code{GeneralMarkovValueEquilibrium} provides
$\theta_s,\varepsilon_{d,s},S(\varepsilon,s),\mu_s$, and a transition matrix $G(t\mid s)$,
with positive tightness, admissible cutoffs, nonnegative transition objects, and
$\sum_tG(t\mid s)=1$.

\decl{GeneralMarkovValueEquilibrium.free\_entry --- paper equation (32)}
\begin{equation}
\mworker{\theta_s}(1-\beta_s)S(\varepsilon_{u,s},s)=c_s. \tag{32}
\end{equation}

\decl{GeneralMarkovValueEquilibrium.cutoff\_zero --- paper equation (33)}
\begin{equation}
S(\varepsilon_{d,s},s)=0. \tag{33}
\end{equation}

\decl{GeneralMarkovValueEquilibrium.surplus\_bellman --- paper equation (34)}
\begin{align}
(r_s+\lambda_s+\mu_s)S(\varepsilon,s)
={}&p_s+\sigma_s\varepsilon-b_s
-\beta_s\mworker{\theta_s}S(\varepsilon_{u,s},s) \notag\\
&+\lambda_s\int\max\{S(x,s),0\}\,dF_s(x) \notag\\
&+\mu_s\sum_tG(t\mid s)\max\{S(\varepsilon,t),0\}. \tag{34}
\end{align}

The separate \code{MarkovEquilibrium} structure is weaker: it contains a reduced
steady-state equilibrium in each state and a row-stochastic $G$, but does not itself
impose equations (32)--(34).

\section{Simulation and accounting definitions}

\source{Definitions.lean and SimulationPath in Equilibrium.lean}

\decl{Primitives.periodCreation --- paper equation (36)}
\begin{equation}
C_t:=a_t(1-N_t). \tag{36}
\end{equation}

\decl{Primitives.immediateScrapping and periodDestruction --- paper equation (37)}
\begin{equation}
D_t^{\rm immediate}:=\int_{\epsd^{\rm old}}^{\epsd^{\rm new}}n(x)\,dx,
\qquad
D_t:=D_t^{\rm immediate}+D_t^{\rm ongoing}. \tag{37}
\end{equation}

\decl{Primitives.nextEmployment --- paper equation (38)}
\begin{equation}
N_{t+1}:=N_t+C_t-D_t. \tag{38}
\end{equation}

\code{SimulationPath P} contains sequences $N_t,a_t,D_t^{\rm immediate},
D_t^{\rm ongoing},C_t,D_t$ and requires (36)--(38) at every natural-numbered
period. The structure does not impose a particular formula on
$D_t^{\rm ongoing}$; that sequence is supplied by an instance.

\section{Kernel-checked theorem catalogue}
\label{sec:theorems}

This section translates every theorem named in \leanfile{Results.lean}. Each row is
also checked by an \code{assert\_no\_sorry} command in \leanfile{Verification.lean}.

\subsection{Value equations and cutoff algebra}

\begin{longtable}{@{}p{0.31\linewidth}p{0.64\linewidth}@{}}
\toprule
\textbf{Lean theorem} & \textbf{Exact mathematical content} \\
\midrule
\endhead
\breakcode{surplus_bellman} & Every \code{ValueEquilibrium} satisfies (8). The proof
adds (5) and (6), subtracts (7), and uses (3)--(4). \\
\breakcode{J_cutoff_eq_zero} & Every \code{ValueEquilibrium} satisfies
$J(\epsd)=0$, using $S(\epsd)=0$, (3), and (4). \\
\breakcode{vacancy_free_entry_value} & Under $r>0$, equations (1)--(2) imply
$q(\theta)J(\epsu)=c$. \\
\breakcode{S_cutoff_self} & $S_d(\epsd;\epsd)=0$. \\
\breakcode{surplus_difference} &
$S_d(\varepsilon_2;\epsd)-S_d(\varepsilon_1;\epsd)
=\sigma(\varepsilon_2-\varepsilon_1)/(r+\lambda)$. \\
\breakcode{nash_shares_exhaust_surplus} &
$J_d(\varepsilon;\epsd)+[W-U]_d(\varepsilon;\epsd)=S_d(\varepsilon;\epsd)$. \\
\breakcode{surplus_nonneg} & Under the maintained sign assumptions,
$\epsd\le\varepsilon\Rightarrow S_d(\varepsilon;\epsd)\ge0$. \\
\breakcode{J_cutoff_nonneg} & Under the same assumptions,
$\epsd\le\varepsilon\Rightarrow J_d(\varepsilon;\epsd)\ge0$. \\
\bottomrule
\end{longtable}

\subsection{Matching, job creation, and destruction}

\begin{longtable}{@{}p{0.35\linewidth}p{0.60\linewidth}@{}}
\toprule
\textbf{Lean theorem} & \textbf{Exact mathematical content} \\
\midrule
\endhead
\breakcode{destructionHazard_nonneg} & $\lambda F(\epsd)\ge0$. \\
\breakcode{destructionHazard_mono} & $\epsd\mapsto\lambda F(\epsd)$ is
nondecreasing. \\
\breakcode{mWorker_pos} & $\theta>0\Rightarrow\mworker{\theta}>0$. \\
\breakcode{jobCreation_iff_freeEntryValue} & If $\epsd<\epsu$, equation (13) is
equivalent to $q(\theta)J_d(\epsu;\epsd)=c$. \\
\breakcode{jobCreationCurve_slopes_down} & For two points satisfying (13), if
$\varepsilon_1<\varepsilon_2<\epsu$, then $\theta_2<\theta_1$ (using that $q$ is
nonincreasing). \\
\breakcode{dispersionCutoffResponse_pos} & If
$b+[\beta c/(1-\beta)]\theta<p$, then $R_\sigma(\theta,\epsd)>0$. This proves the
sign of the defined right-hand side (11R), not an implicit-function theorem. \\
\breakcode{boomJobCreation_zero_transition_iff} & Setting $\mu=0$ in (30) is
equivalent to (13), with the boom cutoff. \\
\breakcode{anticipation_reduces_boom_tightness} & Fix cutoffs with
$\epsd^*<\epsd<\epsu$ and $\mu>0$. If the denominator of (30) is positive, a
no-transition point satisfying (13) and an anticipated point satisfying (30) obey
$\theta^*_{\mu}<\theta^*_0$. \\
\bottomrule
\end{longtable}

\subsection{Flows and stationary accounting}

\begin{longtable}{@{}p{0.35\linewidth}p{0.60\linewidth}@{}}
\toprule
\textbf{Lean theorem} & \textbf{Exact mathematical content} \\
\midrule
\endhead
\breakcode{beveridge_denominator_pos} & If $\theta>0$, then
$\lambda F(\epsd)+\mworker{\theta}>0$. \\
\breakcode{creationFlow_mono} & Holding $u\ge0$ fixed, creation is nondecreasing in
the worker meeting rate. \\
\breakcode{destructionFlow_mono} & Holding $u\le1$ fixed, destruction is
nondecreasing in $\epsd$. \\
\breakcode{aggregateShock_flow_orders} & Conditional on
$\theta_0\le\theta_1$ and $\varepsilon_{d,1}\le\varepsilon_{d,0}$, creation weakly rises and destruction
weakly falls. The equilibrium ordering is a hypothesis. \\
\breakcode{dispersionShock_flow_orders} & Conditional on
$\theta_0\le\theta_1$ and $\varepsilon_{d,0}\le\varepsilon_{d,1}$, creation and destruction both weakly
rise. The equilibrium ordering is a hypothesis. \\
\breakcode{gross_flows_balance} & Every \code{SteadyStateEquilibrium} satisfying the
maintained assumptions has $C=D$. \\
\breakcode{unemploymentDrift_eq_zero} & Every such steady state has $\dot u=0$. \\
\code{v\_eq\_}$\theta$\code{\_mul\_u} & The equilibrium vacancy definition gives
$v=\theta u$. \\
\bottomrule
\end{longtable}

\subsection{Qualitative co-movement and period accounting}

Lean defines
\begin{align*}
\operatorname{MovesOpposite}(x_0,x_1,y_0,y_1)
&:\Longleftrightarrow (x_1-x_0)(y_1-y_0)<0,\\
\operatorname{MovesTogether}(x_0,x_1,y_0,y_1)
&:\Longleftrightarrow 0<(x_1-x_0)(y_1-y_0).
\end{align*}

\begin{longtable}{@{}p{0.38\linewidth}p{0.57\linewidth}@{}}
\toprule
\textbf{Lean theorem} & \textbf{Exact mathematical content} \\
\midrule
\endhead
\breakcode{aggregateShock_negative_comovement} & From strict hypotheses
$C_0<C_1$ and $D_1<D_0$, it follows that creation and destruction satisfy
\code{MovesOpposite}. \\
\breakcode{dispersionShock_positive_comovement} & From strict hypotheses
$C_0<C_1$ and $D_0<D_1$, it follows that the flows satisfy \code{MovesTogether}. \\
\breakcode{unemploymentDrift_pos_iff} &
$\dot u>0\Longleftrightarrow C<D$. \\
\breakcode{employment_change_identity} &
$(N+C-D)-N=C-D$. \\
\breakcode{downturn_destruction_exceeds_ongoing} & If
$D^{\rm immediate}>0$, then
$D^{\rm ongoing}<D^{\rm immediate}+D^{\rm ongoing}$. \\
\breakcode{boom_without_scrapping} &
$0+D^{\rm ongoing}=D^{\rm ongoing}$. \\
\breakcode{SimulationPath.employment_change} & Every simulation path satisfies
$N_{t+1}-N_t=C_t-D_t$. \\
\bottomrule
\end{longtable}

\section{What is proved, assumed, or only represented}
\label{sec:scope}

\begin{longtable}{@{}p{0.20\linewidth}p{0.27\linewidth}p{0.46\linewidth}@{}}
\toprule
\textbf{Status} & \textbf{Topic} & \textbf{Current formal content} \\
\midrule
\endhead
\status{green}{PROVED} & Value-equation aggregation & Equations (3)--(7) imply (8);
free entry implies the zero-profit value equation; cutoff firm value is zero. \\
\status{green}{PROVED} & Stationary accounting & The Beveridge equation implies
$C=D$ and $\dot u=0$; period employment changes by $C-D$. \\
\status{green}{PROVED} & Conditional monotonicity & Hazard, creation, destruction,
JC-locus, and anticipation statements follow from the explicitly listed sign and
ordering hypotheses. \\
\status{amber}{ASSUMED} & Aggregate-shock equilibrium ordering & The rise in
$\theta$ and fall in $\epsd$ are fields or hypotheses before the flow-order and
co-movement lemmas are applied. \\
\status{amber}{ASSUMED} & Dispersion-shock equilibrium ordering & The joint rise in
$\theta$ and $\epsd$ is a hypothesis of the corresponding flow theorem. The code
proves positivity of (11R) under a net-price condition, but not the full joint
implicit comparative static. \\
\status{amber}{REPRESENTED} & Markov equilibrium & Equations (32)--(34) are encoded as
fields of a value-equilibrium structure; no existence, uniqueness, or solution theorem
is currently proved. \\
\status{amber}{REPRESENTED} & Simulation & Equations (36)--(38) and the immediate
scrapping integral are encoded. Calibration, numerical solution, and the paper's
simulated series are not reproduced. \\
\status{red}{OUT OF SCOPE} & Full paper replication & The current Lean development
does not formalize every equation, derive all cyclical orderings from primitives,
prove equilibrium existence/uniqueness, or replicate the Section 5 calibration. \\
\bottomrule
\end{longtable}

\newpage
\thispagestyle{fancy}
This classification follows the actual types in the Lean files. It is the key manual
audit: an implication whose economic conclusion is already a hypothesis is logically
valid but does not by itself reproduce the paper's derivation of that hypothesis.

\section{Verification ledger}

\subsection{No-sorry audit}

\leanfile{Verification.lean} imports \leanfile{Results.lean} and applies
\code{assert\_no\_sorry} to all 31 declarations listed in
Section~\ref{sec:theorems}. The command asks Lean to collect the axioms used by each
declaration and fails if \code{sorryAx} occurs. Therefore a successful dedicated
build establishes both elaboration and absence of unfinished \code{sorry} proofs for
this theorem set.

\subsection{Reproduction commands}

From the repository root:
\begin{verbatim}
lake build Lean4Tutorial.i002_replicate_MP
latexmk -pdf -interaction=nonstopmode -halt-on-error \
  -output-directory=output/pdf \
  output/pdf/mp1994_lean_formalization.tex
\end{verbatim}

\subsection{Manual comparison checklist}

\begin{enumerate}
  \item Compare the primitive definitions and signs in Sections 2--3 with the paper.
  \item Compare equations (1)--(15), (20), (24), (29)--(30), (32)--(34), and
        (36)--(38), paying special attention to integration measures and denominators.
  \item Check whether the paper's $m(\theta,1)$ convention matches
        $\theta q(\theta)$ throughout.
  \item Distinguish theorems that derive an ordering from those that accept it as a
        hypothesis or structure field.
  \item Decide whether the omitted existence, uniqueness, implicit-function, and
        calibration arguments are required for the intended meaning of ``replication.''
\end{enumerate}

\section*{Reference}
\addcontentsline{toc}{section}{Reference}

Dale T. Mortensen and Christopher A. Pissarides (1994), ``Job Creation and Job
Destruction in the Theory of Unemployment,'' \emph{Review of Economic Studies},
61(3), 397--415. The locally supplied paper is
\leanfile{Mortensen-JobCreationJob-1994.pdf}.

\end{document}
